\section{Conclusion}
\label{sec:conclusion}

\policy{} removes padded OT work from batches whose problems require different
numbers of iterations.
It starts from standard fixed-width execution of independent Sinkhorn problems;
the selected solver's stopping rule determines completion, and dynamic compaction narrows every subsequent per-problem kernel. The survival-curve
identity counts the problem updates that disappear. The opportunity--conversion
decomposition then maps survivor width through the measured backend cost curve
and subtracts control overhead, giving the exact matched-path break-even
condition.

Matched fixed-width/dynamic-compaction comparisons deliver
\CMetroCorrectedEndpointSpeedup{} on the four-worker MetroPT-3 endpoint,
$1.054\times$ on the ImageNet-32 training solver field, and $1.616$--$3.110\times$
in the registered MetroPT scale cells. A non-cold setting that uses one
initialization for multiple problems gives
$1.421\times$ with 30.62\% fewer problem-update slots. Independent OTT-JAX and
PyKeOps implementations remove completed problems through independent backends.

Matched controls identify the mechanism. Batched checking and the
Sinkhorn-specific screen reduce inspection cost; logical masking leaves the kernel width unchanged; dynamic compaction removes OT work. In the registered
MetroPT and Packer19 endpoint windows, dynamic compaction also reduces measured GPU
energy.
Grouping while holding total iteration counts fixed shows that greater
variation in the required iteration numbers within a window creates more
removable padding. Direct A100 width sweeps show a wave-shaped conversion curve:
small widths can share one cost level, while larger problems respond to the
width reduction earlier.
width reduction earlier. The reported applications pair speed with
achieved residuals and consumer or training endpoints. Fast inner kernels make
each update cheaper; DrainSinkhorn stops issuing
those updates to completed problems.

\paragraph{Reproducibility.}
The release artifact records source and input identities, commands, software and
GPU environments, raw outputs, numerical checks, and paired analyses.
